\chapter{致谢与参考}
\label{front:acknowledgements}
\vspace{-4mm}
\noindent{\color{CoverTeal}\rule{\linewidth}{1.2pt}}

\noindent\textbf{\large\color{Teal}感谢}

感谢下列著作、讲义与文章的作者，感谢认真讲授数学、分享解法的老师和同学，也感谢开放教材、公开课程及数学工具的建设者。愿本册帮助读者把直觉写成推导，把技巧连成体系。欢迎读者交流解法、提出勘误。

\smallskip
\noindent\textbf{\large\color{Teal}参考文献与延伸阅读}
\par\smallskip
{\setlength{\columnsep}{7mm}\setlength{\multicolsep}{0pt}
\begin{multicols}{2}
\fontsize{9.3}{12.5}\selectfont
\setlength{\parskip}{0pt}\interlinepenalty=10000
\begin{enumerate}[label=\textcolor{Teal}{[\arabic*]},leftmargin=2.2em,itemsep=2pt,topsep=1pt,parsep=0pt]
\item 虚调子：《不定积分攻略》（v4.0，封面题名《不定积分之潮》）；《不定积分之双元法》。
\item 唐绍东（TangShaodong）：《微积分笔记》（特别版，2024年9月8日）。
\item 金玉明、顾新身、毛瑞庭：《积分的方法与技巧》，中国科学技术大学出版社，2017年。
\item 孙清华、李金兰、孙昊：《常微分方程内容、方法与技巧》。
\item 夜雨（主讲）：《微分方程笔记》；《级数概念与选择题》。
\item 作者未署名：《极限求法大总结》（第三版）。
\item 作者未署名：《高等数学笔记》。
\item Steven J. Kifowit：\textit{More Proofs of Divergence of the Harmonic Series}（调和级数发散证明）。
\item Juan Arias de Reyna：\textit{Pointwise Convergence of Fourier Series}（傅里叶级数的点收敛），Springer，2002年。
\item 杨圣全、石明：《级数求和的八种方法》，《高等数学研究》，2013年。
\item 李绍刚、徐安农：《二阶常系数线性微分方程特解的微分算子法》，《桂林电子科技大学学报》，2008年。
\item 胡昊宇：《无穷级数的求和方法》，《教育现代化》，2019年。
\item 周杰荣：《正项级数敛散性的对数判别法与拉贝判别法》。
\item 静流：《半天速通大学生数学竞赛——级数计算》；《非数一元积分突击》；《三天速成大学生数学竞赛——多元积分计算》；《速通大学生数学竞赛（非数）——极限计算》；《速通大学生数学竞赛——微分方程计算》。
\item epsilon-delta（讲义署名）：《中值定理专题讲义》；《一元函数微分学》；《数列极限专题》；《保号性专题》；《不定积分与定积分专题》；《反常积分专题》；《方程根与函数零点问题专题》；《泰勒公式专题》；《凹凸性专题》；《切比雪夫不等式补充三题》；《积分不等式专题》；《多元函数微分学》；《二重积分专题》；《微积分应用》；《级数专题》；《三重积分》；《曲线曲面积分专题》。
\item Gilbert Strang、Edwin Herman：\href{https://openstax.org/books/calculus-volume-1/pages/preface}{\textit{Calculus}}，OpenStax。
\item MIT OpenCourseWare：\href{https://ocw.mit.edu/courses/18-03-differential-equations-spring-2010/}{\textit{18.03 Differential Equations}}，2010年春季。
\item NIST：\href{https://dlmf.nist.gov/}{\textit{Digital Library of Mathematical Functions}}，重点参照第9、10、19章。
\item 宾夕法尼亚大学分析讲义：\href{https://www2.math.upenn.edu/~gressman/analysis/14-plancherel.html}{\textit{The Fourier Transform on L²}}。
\end{enumerate}
\end{multicols}}

\smallskip
\noindent\textbf{\large\color{Teal}联系与勘误}

\begin{tcolorbox}[colback=Teal!3,colframe=Teal!35,boxrule=.4pt,before skip=5pt,after skip=4pt]
\small
\textbf{作者}\quad 雪落长安无声\hfill
\textbf{知乎}\quad \href{https://www.zhihu.com/people/black-despair}{zhihu.com/people/black-despair}\par\smallskip
\textbf{QQ邮箱}\quad \href{mailto:3547908579@qq.com}{3547908579@qq.com}\quad （勘误请附章节、题号或公式位置）
\end{tcolorbox}

\vfill
{\footnotesize\color{black!60}封面摄影：freestocks.org，\href{https://commons.wikimedia.org/wiki/File:Open_book_(26421751605).jpg}{\textit{Open book}}，CC0 1.0。\par}
